\documentclass[12pt,a4paper]{article}

% --- Pacotes de Suporte ---
\usepackage{cmap} % corrige o mapeamento Unicode de ligaduras (fi/ff/fl)
\usepackage[utf8]{inputenc}
\usepackage[T1]{fontenc}
\usepackage[expansion=false]{microtype}
\usepackage[brazilian]{babel}
\usepackage{amsmath}
\usepackage{amsfonts}
\usepackage{amssymb}
\usepackage{geometry}
\usepackage{bm}
\usepackage{graphicx}
\usepackage{float}
\usepackage{booktabs}
\usepackage{hyperref}
\usepackage{fancyhdr}
\usepackage{titlesec}
\usepackage{tabularx}
\usepackage{longtable}
\usepackage{tikz}
\usetikzlibrary{patterns}
\usepackage{pgfplots}
\pgfplotsset{compat=1.18}

\hyphenation{
  off-shore Off-shore
  con-tí-nuo con-tí-nu-a
  di-fe-ren-ci-al di-fe-ren-ci-ais
  lo-ga-rit-mo lo-ga-rit-mos
  tri-go-no-mét-ri-ca tri-go-no-mét-ri-cas
  a-pro-xi-ma-ção a-pro-xi-ma-ções
}

\numberwithin{equation}{section}
\numberwithin{table}{section}
\numberwithin{figure}{section}

\geometry{margin=2.5cm, headheight=15pt}
\addtolength{\topmargin}{-0.5pt}
\setlength{\emergencystretch}{1em}

\hypersetup{
	colorlinks=true,
	linkcolor=black,
	citecolor=black,
	urlcolor=black,
	pdftitle={Notações e Cálculo Básico},
	pdfauthor={Marcelo Moreira Mejias},
	pdfsubject={Pré-requisitos para Máquinas Rotodinâmicas},
	pdfkeywords={notação matemática, cálculo, vetores, equações diferenciais, aproximações}
}

\pagestyle{fancy}
\fancyhf{}
\renewcommand{\sectionmark}[1]{\markboth{#1}{}}
\fancyhead[L]{\nouppercase{\leftmark}}
\fancyhead[R]{\thepage}
\renewcommand{\headrulewidth}{0.4pt}

% --- Identificação de volume/revisão ---
\newcommand{\colecaotitulo}{Enciclopédia de Máquinas Rotodinâmicas}
\newcommand{\volumeatual}{Volume 0}
\newcommand{\revdata}{Rev. 12/09/2026}

% --- Identidade visual dos capítulos ---
\newcommand{\chapterword}{Capítulo}
\titleformat{\section}
	{\normalfont\Large\bfseries\raggedright}
	{\chapterword~\thesection}{1em}{}
	[\vspace{-0.5em}\rule{\textwidth}{0.4pt}]


\begin{document}


	% --- CAPA ---
	\begin{titlepage}
		\centering
		\vspace*{1.5cm}
		{\large \textsc{\colecaotitulo\ --- \volumeatual}} \\[1cm]
		\rule{\textwidth}{0.4pt} \\[0.8cm]
		{\Huge \textbf{Notações e Cálculo Básico}} \\[0.5cm]
		{\Large \textbf{Pré-requisitos para o Estudo de Máquinas Rotodinâmicas}} \\[0.8cm]
		\rule{\textwidth}{0.4pt} \\[1.5cm]

		% Placeholder para ilustração conceitual do Volume 0
		\vspace{5cm}

		\vfill
		{\large \textbf{Autor: Engenheiro Marcelo Moreira Mejias, M.Sc.}} \\[0.5cm]

		{\large \today} \\[0.3cm]
		{\normalsize \revdata}
	\end{titlepage}

	% --- DEDICATÓRIA ---
	\newpage
	\thispagestyle{empty}
	\vspace*{\fill}
	\begin{flushright}
		\textit{Aos que reaprendem o básico sem vergonha,}
		\textit{e aos que nunca o esqueceram.}
	\end{flushright}
	\vspace*{\fill}


% --- PREFÁCIO ---
\newpage
\section*{Prefácio}
\addcontentsline{toc}{section}{Prefácio}

Esta coleção pede do leitor um repertório mínimo: ler e manipular símbolos
matemáticos, operar com unidades, derivar, integrar e resolver as equações
diferenciais mais simples. Quem chega da formação recente talvez precise
apenas revisar; quem está há anos fora da sala de aula, ou quem nunca viu o
formalismo, precisa de um lugar para começar. O Volume~0 é esse lugar.

A intenção não é substituir os livros de cálculo e de física básica, mas
reunir, num único percurso curto, exatamente o que os Volumes~I e~II usam ---
nem mais, nem menos. Cada capítulo declara o que apresenta, mostra o que é
operacional e fecha com uma síntese que aponta onde a ferramenta reaparece
nos volumes seguintes. O leitor que domina o assunto de um capítulo pode
pulá-lo sem prejuízo; o leitor que duvida encontra aqui o mínimo necessário,
sem rodeios.

Nussenzveig~\cite{nussenzveig} cobre a física básica que antecede esta
notação; Guidorizzi~\cite{guidorizzi2013} cobre o cálculo em profundidade.
Para o aprofundamento, as referências de classe mundial são as de
Feynman~\cite{feynman1963} na física e de Apostol~\cite{apostol1967} e
Spivak~\cite{spivak2008} no cálculo. Este volume é o caminho curto entre
essas obras e a engenharia de máquinas.

% --- SUMÁRIO ---
\newpage
\tableofcontents
\newpage


%=============================================================================
% CAPÍTULO 1 — INTRODUÇÃO AO VOLUME 0
%=============================================================================
\section{Introdução ao Volume 0}\label{sec:v0-intro}

\subsection{Para quem é este volume}

Há dois leitores em mente. O primeiro é quem retoma os estudos depois de
anos de prática e percebe que o vocabulário simbólico esqueceu mais rápido
que o juízo físico --- sabe o que o problema pede, mas não lê mais a equação
com fluidez. O segundo é quem começa do zero relativo: tem firmeza em
aritmética e no raciocínio de causa e efeito, mas nunca sistematizou
funções, derivadas e integrais. Este volume serve aos dois. Para o primeiro,
é um relembrar ordenado; para o segundo, um aprender enxuto, com tudo o que
será cobrado depois --- e nada do que não será.

\subsection{O que este volume cobre --- e o que não cobre}

O recorte é definido pelo uso, não pela tradição. Estão aqui a notação
simbólica, as unidades e a análise dimensional, as funções elementares, os
vetores e o essencial de álgebra linear, o cálculo diferencial e integral em
uma e várias variáveis, as equações diferenciais ordinárias de primeira e
segunda ordem e as séries de aproximação. Não estão aqui a teoria dos
números reais, as demonstrações de existência e unicidade, a integração
avancada de funções especiais nem a teoria de espaços vetoriais: nada disso
aparece nos Volumes~I e~II, e portanto não tem lugar neste volume.

\subsection{Como usar este volume}

Os capítulos foram escritos para serem lidos em sequência, mas cada um pode
ser consultado isoladamente: todos usam a mesma notação, estabelecida no
Capítulo~\ref{sec:v0-notacao} e mantida sem exceção até o Volume~II. A
Tabela~\ref{tab:v0-c1-1} resume o papel de cada capítulo e o ponto dos
volumes seguintes em que ele reaparece.

\begin{table}[H]
\centering
\caption[Mapa do Volume 0 e suas pontes com os volumes I e II.]{%
Mapa do Volume 0: cada capítulo ensina uma ferramenta e indica onde ela é
usada nos Volumes I e II desta coleção.}
\label{tab:v0-c1-1}
\begin{tabularx}{\textwidth}{lXl}
\toprule
\textbf{Cap.} & \textbf{Ferramenta} & \textbf{Onde reaparece} \\
\midrule
\ref{sec:v0-notacao} & Notação e linguagem simbólica & Toda a coleção \\
\ref{sec:v0-unidades} & Unidades e análise dimensional & Toda a coleção \\
\ref{sec:v0-funcoes} & Funções elementares e gráficos & Caps. 2 a 9 do Vol.~I \\
\ref{sec:v0-vetores} & Vetores e álgebra linear & Caps. 2, 4 e 7 do Vol.~I \\
\ref{sec:v0-derivadas} & Cálculo diferencial & Leis de taxa, Vol.~I \\
\ref{sec:v0-integrais} & Cálculo integral & Balanços, Vol.~I \\
\ref{sec:v0-edo} & Equações diferenciais & Modelagem dinâmica, Vol.~II \\
\ref{sec:v0-series} & Séries e aproximações & Linearização e erros, Vols.~I e~II \\
\bottomrule
\end{tabularx}
\end{table}

Quem já domina um capítulo avança sem culpa; quem duvida resolve os exemplos
à mão --- nenhum passo deste volume exige computador.


%=============================================================================
% CAPÍTULO 2 — NOTAÇÃO E LINGUAGEM SIMBÓLICA
%=============================================================================
\section{Notação e Linguagem Simbólica}\label{sec:v0-notacao}

\subsection{Por que a notação importa}

A notação é a gramática da física aplicada. Uma equação como
$\tau = J_p\,\alpha$ só pode ser lida de cor porque cada símbolo tem papel
fixo: letras em itálico representam grandezas que variam, letras em romano
representam nomes e funções, operadores têm formas próprias. Quem internaliza
essa gramática lê equações como frases; quem não internaliza decora
desenhos. Este capítulo fixa a gramática usada em toda a coleção.

\subsection{Fontes: itálico para variáveis, romano para nomes}

A convenção seguida em todos os volumes é a dos textos clássicos de
matemática aplicada:

\begin{itemize}
\item \textbf{Itálico} para variáveis e grandezas: $x$, $T$ (temperatura),
$p$ (pressão), $t$ (tempo), $\omega$ (velocidade angular).
\item \textbf{Romano} para operadores, funções com nome e unidades:
$\mathrm{d}x$ (diferencial), $\mathrm{sen}$, $\mathrm{ln}$, $\mathrm{exp}$,
$\max$, e as unidades $\mathrm{m}$, $\mathrm{s}$, $\mathrm{Pa}$.
\item \textbf{Negro} para vetores e matrizes: $\bm{v}$, $\bm{F}$,
$\bm{\nabla}T$.
\end{itemize}

O diferencial escreve-se $\mathrm{d}$ romano para distinguir o operador da
variável: $\mathrm{d}t$ lê-se ``um pedacinho de $t$''. O número
$\mathrm{e} \approx 2{,}71828$ (base dos logaritmos naturais) e a unidade
imaginária $\mathrm{i}$ seguem a mesma regra.

\subsection{Operadores e símbolos de uso geral}

A Tabela~\ref{tab:v0-c2-1} reúne os operadores que aparecem em toda a
coleção. Nenhum deles é opcional: todos serão usados nos volumes seguintes.

\begin{table}[H]
\centering
\caption[Operadores matemáticos e seu significado.]{%
Operadores matemáticos usados em toda a coleção.}
\label{tab:v0-c2-1}
\begin{tabularx}{\textwidth}{lX}
\toprule
\textbf{Símbolo} & \textbf{Significado} \\
\midrule
$=$ & igualdade exata \\
$\approx$ & aproximadamente igual (dentro de uma tolerância declarada) \\
$\equiv$ & identidade, definição ou equivalência por hipótese \\
$:= $ & ``é definido como'' (a igualdade aponta da definição para o nome) \\
$\propto$ & proporcional a (igual a menos uma constante) \\
$\rightarrow$, $\to$ & tende a; ou mapeamento de função \\
$\sum$, $\prod$ & somatório e produtório de uma sequência \\
$\mathrm{d}/\mathrm{d}t$ & derivada em relação ao tempo (taxa instantânea) \\
$\partial/\partial x$ & derivada parcial em relação a $x$ \\
$\nabla$ & gradiente (vetor das derivadas parciais) \\
$\int_a^b$ & integral definida entre $a$ e $b$ \\
$\Delta$ & variação finita (ex.: $\Delta T = T_2 - T_1$) \\
$\delta$ & variação infinitesimal ou incerteza, conforme o contexto \\
$\lfloor\cdot\rfloor$, $\lceil\cdot\rceil$ & piso e teto \\
$\forall$, $\exists$ & ``para todo'', ``existe'' \\
\bottomrule
\end{tabularx}
\end{table}

Duas distinções valem reforço. A variação finita $\Delta x = x_2 - x_1$ é um
número comum entre dois estados; o diferencial $\mathrm{d}x$ é um pedacinho
de variável dentro de um processo de limite --- a integral é feita deles. E
a distinção entre $\approx$ e $\equiv$: a primeira admite erro, que deve ser
declarado ou estimável; a segunda não admite erro nenhum, apenas convenção.

\subsection{Letras gregas}

A engenharia de máquinas usa intensivamente as letras gregas, e a mesma
letra costuma ter uso dominante em cada área. A Tabela~\ref{tab:v0-c2-2}
lista as que aparecem nos Volumes~I e~II com seu papel principal.

\begin{table}[H]
\centering
\caption[Letras gregas e seu uso dominante nesta coleção.]{%
Letras gregas mais usadas na coleção e seu significado dominante.}
\label{tab:v0-c2-2}
\begin{tabularx}{\textwidth}{llX}
\toprule
\textbf{Min.} & \textbf{Mai.} & \textbf{Uso dominante} \\
\midrule
$\alpha$ & & aceleração angular; coeficiente de expansão \\
$\beta$ & & razão de compressão volumétrica \\
$\gamma$ & $\Gamma$ & razão de calores específicos $c_p/c_v$ \\
$\delta$ & $\Delta$ & variação finita; folga (Volume II) \\
$\varepsilon$ & & deformação; eficiência alternativa; rugosidade \\
$\zeta$ & & fator de amortecimento \\
$\eta$ & & eficiência \\
$\theta$ & $\Theta$ & ângulo; temperatura reduzida \\
$\lambda$ & $\Lambda$ & comprimento de onda; autovalor \\
$\mu$ & & viscosidade dinâmica; coeficiente de atrito \\
$\nu$ & & viscosidade cinemática; frequência natural \\
$\rho$ & & massa específica \\
$\sigma$ & $\Sigma$ & tensão; somatório \\
$\tau$ & & torque; tensão de cisalhamento \\
$\phi$ & $\Phi$ & razão de fluxo; função \\
$\psi$ & $\Psi$ & função de corrente (fluidos) \\
$\omega$ & $\Omega$ & velocidade angular; resistência elétrica \\
\bottomrule
\end{tabularx}
\end{table}

Quando uma letra servir a dois usos, o texto declara qual vale na ocasião;
o Volume~I adota essa regra desde seu primeiro capítulo.

\subsection{Índices, expoentes e adornos}

Grandeza nenhuma anda sozinha: a coleção escreve $p_1$ e $p_2$ para pressões
em dois pontos, $T_{h,i}$ para temperatura de entrada do fluido quente, e
$\dot{m}$ (ponto por cima) para vazão mássica --- o ponto codifica
``por unidade de tempo''. Adornos têm significado fixo: linha por cima
($\bar{x}$) para média, til ($\tilde{x}$) para valor adimensionalizado,
chapéu ($\hat{x}$) para estimativa. Subscritos romanos nomeiam
($c_p$ = calor específico \emph{a pressão constante}); subscritos itálicos
indexam ($x_i$, o $i$-ésimo valor).

\subsection{Síntese: premissas e ligação com os volumes seguintes}

Este capítulo não postulou física nenhuma: fixou apenas a escrita. As
convenções aqui estabelecidas --- fontes, operadores, letras gregas,
adornos --- valem sem exceção do Capítulo~\ref{sec:v0-unidades} deste volume
até o capítulo final do Volume~II. Quando o Volume~I escreve
$\bm{\nabla}T$, é o gradiente desta Tabela~\ref{tab:v0-c2-1}; quando escreve
$\eta$, é a eficiência desta Tabela~\ref{tab:v0-c2-2}.


%=============================================================================
% CAPÍTULO 3 — UNIDADES, DIMENSÕES E ORDEM DE GRANDEZA
%=============================================================================
\section{Grandezas, Unidades e Análise Dimensional}\label{sec:v0-unidades}

\subsection{Grandezas e unidades: o número e o rótulo}

Toda grandeza física é um número multiplicado por uma unidade:
``$p = 7{,}0\ \mathrm{bar}$'' significa $7{,}0$ vezes a unidade bar.
O número isolado não significa nada físico --- 7,0 de quê? --- e a unidade
isolada tampouco. A primeira disciplina deste capítulo é nunca separar os
dois: escrever a unidade junto a cada valor, em todos os passos de todos os
cálculos.

As unidades usadas em toda a coleção são as do Sistema Internacional (SI),
resumidas na Tabela~\ref{tab:v0-c3-1} com as grandezas de interesse desta
obra. Unidades fora do SI aparecem apenas quando a prática industrial as
impõe --- o bar para pressão e o rpm para rotação ---, e sempre com conversão
explícita na primeira ocorrência.

\begin{table}[H]
\centering
\caption[Unidades do SI usadas nesta coleção.]{%
Grandezas e unidades do SI usadas nesta coleção.}
\label{tab:v0-c3-1}
\begin{tabularx}{\textwidth}{llX}
\toprule
\textbf{Grandeza} & \textbf{Símbolo} & \textbf{Unidade SI (abrev.)} \\
\midrule
comprimento & $L$ & metro (m) \\
massa & $m$ & quilograma (kg) \\
tempo & $t$ & segundo (s) \\
corrente elétrica & $I$ & ampère (A) \\
temperatura & $T$ & kelvin (K) \\
\midrule
força & $F$ & newton (N = kg\,m\,s$^{-2}$) \\
pressão & $p$ & pascal (Pa = N\,m$^{-2}$) \\
energia, trabalho & $E$, $W$ & joule (J = N\,m) \\
potência & $P$ & watt (W = J\,s$^{-1}$) \\
\bottomrule
\end{tabularx}
\end{table}

Os prefixos do SI multiplicam as unidades por potências de dez. Os que
aparecem na coleção são poucos e devem ser lidos de cor:
$\mathrm{k}$ (quilo, $10^3$), $\mathrm{M}$ (mega, $10^6$),
$\mathrm{G}$ (giga, $10^9$), $\mathrm{m}$ (mili, $10^{-3}$),
$\mu$ (micro, $10^{-6}$). Assim, $1\ \mathrm{kW} = 10^3\ \mathrm{W}$ e
$1\ \mu\mathrm{m} = 10^{-6}\ \mathrm{m}$. Dois cuidados permanentes: o
prefixo $\mathrm{m}$ minúsculo é mili, enquanto $\mathrm{M}$ maiúsculo é
mega --- seis ordens de grandeza de diferença; e prefixos não se combinam
(diz-se $\mathrm{kg}$, nunca ``$\mathrm{k}\mu\mathrm{g}$'').

\subsection{Homogeneidade dimensional}

Uma equação física só pode somar grandezas de mesma dimensão: não se soma
metros a segundos, nem joules a watts. Essa exigência, chamada
\textbf{homogeneidade dimensional}, é o verificador mais barato que existe:
qualquer equação que some termos de dimensões diferentes está errada, antes
de qualquer cálculo numérico. As dimensões fundamentais da mecânica são
comprimento $[L]$, massa $[M]$ e tempo $[t]$; toda grandeza mecânica se
decompõe nelas. A Equação~\eqref{eq:v0-c3-1} mostra a verificação para a
pressão:

\begin{equation}\label{eq:v0-c3-1}
[p] = \frac{[F]}{[A]} = \frac{[M]\,[L]\,[t]^{-2}}{[L]^2}
    = [M]\,[L]^{-1}\,[t]^{-2}.
\end{equation}

A análise dimensional faz mais que detectar erro: ela \emph{descobre} a forma
das relações. Se um período de oscilação depende só de uma massa $m$ e de
uma rigidez $k$ (dimensão $[M][t]^{-2}$), a única combinação com dimensão de
tempo é $\sqrt{m/k}$ --- a constante numérica a análise não dá, mas a
dependência inteira ela dá. O Volume~I usa esse argumento nos grupos
adimensionais do Capítulo~4, e o Volume~II o repete no dimensionamento
preliminar.

\subsection{Notação científica e ordem de grandeza}

A notação científica escreve qualquer número como
$m \times 10^n$, com $1 \le m < 10$ e $n$ inteiro:
$c = 3{,}00 \times 10^8\ \mathrm{m/s}$. Ela existe para dois motivos
operacionais. O primeiro é comparar: saber se um resultado é
$2 \times 10^{-3}$ ou $2 \times 10^{3}$ é saber se a conta está no
universo certo, antes de conferir qualquer dígito. O segundo é o
\textbf{estimativa de ordem de grandeza}: arredondar cada fator para a
potência de dez mais próxima e multiplicar as potências. Uma bomba de
$3{,}2\ \mathrm{kW}$ operando $6{,}5\ \mathrm{h/dia}$ gasta, em ordem de
grandeza, $10^3 \times 10^1 = 10^4\ \mathrm{Wh}$ por dia --- dez ve menos
que mil, dez vezes mais que cem. Quem estima antes de calcular nunca
confunde kW com MW.

\subsection{Algarismos significativos e arredondamento}

Um instrumento que lê $7{,}0\ \mathrm{bar}$ afirma algo mais forte que um
que lê $7\ \mathrm{bar}$: afirma que o valor está entre $6{,}95$ e
$7{,}05$. Os \textbf{algarismos significativos} codificam essa afirmação:
o último dígito escrito é o incerto, todos os anteriores são firmes. As
regras operacionais são poucas:

\begin{itemize}
\item na multiplicação e divisão, o resultado tem tantos significativos
  quanto o fator com menos significativos;
\item na adição e subtração, o resultado tem tantas casas decimais quanto
  o termo com menos casas decimais;
\item o arredondamento é feito no resultado final, nunca nos passos
  intermediários --- guardar dígitos extras durante a conta não é
  desonestidade, é precaução.
\end{itemize}

A regra do resultado final importa mais que as outras duas: uma conta com
dez passos em que cada passo arredonda acumula erro sistemático. O Volume~I
trata a propagação formal de incerteza no capítulo de instrumentação; aqui
basta a disciplina acima.

\subsection{Síntese: premissas e ligação com os volumes seguintes}

Este capítulo postulou apenas que grandezas têm dimensão e que equações
somam dimensões iguais. Com isso, o leitor já pode conferir qualquer
equação dos volumes seguintes: a homogeneidade dimensional da
Equação~\eqref{eq:v0-c3-1} se reaplica em cada lei de taxa do Volume~I, e a
estimativa de ordem de grandeza é o primeiro teste de plausibilidade de
qualquer número de engenharia.


%=============================================================================
% CAPÍTULO 4 — FUNÇÕES ELEMENTARES E GRÁFICOS
%=============================================================================
\section{Funções Elementares e Gráficos}\label{sec:v0-funcoes}

\subsection{O que é uma função}

Uma função associa a cada valor de uma variável independente $x$ um único
valor da variável dependente $y$, escrito $y = f(x)$. A restrição ``único
valor'' não é burocracia: uma curva de pressão de um compressor que para
cada vazão desse duas pressões não é uma função $p(\dot{m})$ --- falta-lhe
informação (em geral, a condição do processo). Em engenharia, especificar a
função exige especificar o \emph{domínio} (para quais $x$ vale) e a
\emph{condição} sob a qual vale; o Volume~I repete esse aviso em cada lei
de taxa.

\subsection{Potências, polinômios e raízes}

A família das potências $f(x) = x^n$ cobre grande parte da geometria
elementar: $n=1$ é a reta, $n=2$ a parábola, $n=3$ a cúbica, e expoentes
negativos dão as leis de decaimento --- a lei de resfriamento de Newton usa
decaimento exponencial, mas a perda de carga em regime laminar
($\Delta p \propto \dot{m}$) e turbulento ($\Delta p \propto \dot{m}^2$)
são potências. Polinômios somam potências inteiras:
$f(x) = a_0 + a_1 x + a_2 x^2 + \dots + a_n x^n$, e a curva
torque--velocidade de um motor de indução é bem aproximada por um
polinômio de segundo grau no seu trecho de operação --- como mostra o
Capítulo~7 do Volume~I.

\subsection{Exponencial e logaritmo}

A função exponencial $f(x) = \mathrm{e}^{x}$ e o logaritmo natural
$f(x) = \ln x$ são inversas uma da outra e aparecem em toda parte onde a
taxa de variação de uma grandeza é proporcional à própria grandeza:
resfriamento, descarga de pressão, crescimento de falha. Suas propriedades
operacionais, Equações~\eqref{eq:v0-c4-1} a~\eqref{eq:v0-c4-3}, devem ser
manipuladas sem hesitação:

\begin{equation}\label{eq:v0-c4-1}
\mathrm{e}^{a+b} = \mathrm{e}^{a}\,\mathrm{e}^{b},
\qquad
\mathrm{e}^{a\,b} = \left(\mathrm{e}^{a}\right)^{b},
\end{equation}
\begin{equation}\label{eq:v0-c4-2}
\ln(a\,b) = \ln a + \ln b,
\qquad
\ln(a^{b}) = b\,\ln a,
\qquad
\ln \mathrm{e}^{x} = x,
\end{equation}
\begin{equation}\label{eq:v0-c4-3}
a^{x} = \mathrm{e}^{x\,\ln a}.
\end{equation}

O logaritmo transforma multiplicação em adição e potência em produto ---
é a calculadora dos tempos sem calculadora, e continua útil porque
lineariza gráficos, como se vê adiante. A exponencial tem ainda a
propriedade $f'(x) = f(x)$: é a única função igual à sua própria taxa de
variação, e por isso governa todos os processos de primeira ordem do
Volume~II.

\subsection{Funções trigonométricas}

As funções $\mathrm{sen}\,\theta$, $\mathrm{cos}\,\theta$ e
$\mathrm{tg}\,\theta$ medem lados do triângulo inscrito no círculo unitário,
Figura~\ref{fig:v0-c4-1}. Toda a teoria de vibrações do Volume~I é escrita
nessas funções, porque todo movimento oscilatório simples é projeção de
rotação.

\begin{figure}[H]
\centering
\begin{tikzpicture}[font=\footnotesize, scale=2.2]
\draw[->] (-1.25,0) -- (1.35,0) node[below left=1pt] {$x$};
\draw[->] (0,-1.25) -- (0,1.35) node[below left=1pt] {$y$};
\draw (0,0) circle (1);
\draw[thick,->] (0,0) -- (35:1);
\draw (0,0) -- ({cos(35)},0);
\draw[dashed] (35:1) -- ({cos(35)},0);
\draw (0.16,0) arc (0:35:0.16) node[midway,right=2pt] {$\theta$};
\draw ({cos(35)},0.035) -- ({cos(35)-0.035},0.035) -- ({cos(35)-0.035},0);
\node[below] at ({cos(35)},0) {$\cos\theta$};
\draw[dashed] (0,{sin(35)}) -- (35:1);
\draw (-0.035,{sin(35)}) -- (-0.035,{sin(35)-0.035}) -- (0,{sin(35)-0.035});
\node[left] at (0,{sin(35)}) {$\mathrm{sen}\,\theta$};
\end{tikzpicture}
\caption[Círculo unitário: seno e cosseno como projeções.]{%
O círculo unitário define seno e cosseno como projeções do raio unitário:
o ponto do círculo a ângulo $\theta$ tem coordenadas
$(\cos\theta,\ \mathrm{sen}\,\theta)$.}
\label{fig:v0-c4-1}
\end{figure}

As identidades que a coleção usa sem aviso são as de soma de arcos e o
teorema fundamental, Equação~\eqref{eq:v0-c4-4}:

\begin{equation}\label{eq:v0-c4-4}
\mathrm{sen}^2\theta + \cos^2\theta = 1,
\qquad
\mathrm{sen}(a \pm b) = \mathrm{sen}\,a\cos b \pm \cos a\,\mathrm{sen}\,b.
\end{equation}

Os argumentos das funções trigonométricas são sempre ângulos em
\emph{radianos} nesta coleção --- a frequência natural $\omega_n$ aparece
multiplicando $t$ dentro do seno, e $\omega t$ só tem dimensão de ângulo
porque o radiano é adimensional.

\subsection{Funções de duas variáveis e superfícies}

Uma função $f(x,y)$ associa um valor a cada par ordenado do plano; seu
gráfico é uma superfície. Na prática desta coleção, funções de duas
variáveis aparecem como \emph{mapas}: a eficiência politrópica como função
da razão de compressão e da rotação, a pressão de descarga como função da
vazão e da abertura de válvula. Ler esses mapas --- fixar uma variável e
ler a curva resultante na outra --- é exatamente o que o cálculo de
derivadas parciais formaliza no Capítulo~\ref{sec:v0-derivadas}.

\subsection{Leitura de gráficos: linear, semilog e log--log}

Gráficos em escala linear mostram somas como deslocamentos; gráficos em
escala logarítmica mostram \emph{potências como retas}. Se
$y = C\,x^n$, então $\ln y = \ln C + n\,\ln x$: num gráfico log--log, a
lei de potência vira uma reta de inclinação $n$. Se
$y = C\,\mathrm{e}^{a x}$, então $\ln y = \ln C + a x$: num gráfico
semilog (só o eixo $y$ em log), a exponencial vira uma reta de inclinação
$a$. Os ensaios do Volume~II usam esses artifícios para ler expoentes de
politropas e constantes de tempo direto do gráfico.

\subsection{Síntese: premissas e ligação com os volumes seguintes}

As funções deste capítulo --- potência, exponencial, logaritmo,
trigonométricas --- são o vocabulário completo dos modelos dos volumes
seguintes: nenhuma outra função aparece nos Volumes~I e~II sem ser
explicitamente construída a partir destas. O círculo unitário da
Figura~\ref{fig:v0-c4-1} retorna nas vibrações mecânicas; as propriedades
da exponencial, nos sistemas de primeira ordem; a linearização log--log,
na leitura de curvas de ensaio.


%=============================================================================
% CAPÍTULO 5 — VETORES E ÁLGEBRA LINEAR ESSENCIAL
%=============================================================================
\section{Vetores e Álgebra Linear Essencial}\label{sec:v0-vetores}

\subsection{Por que vetores}

Força, velocidade, aceleração, gradiente de temperatura e campo magnético
têm direção, sentido e magnitude: tratá-los como números escalares
produz equações incompletas ou erradas. O vetor é a estrutura matemática
que registra as três informações de uma vez. Nesta coleção, vetores
escreve-se em negro ($\bm{v}$) e seus módulos em itálico ($v$).

\subsection{O que é um vetor}

Um vetor $\bm{v}$ num espaço tridimensional é uma terna ordenada de
componentes numa base escolhida, Equação~\eqref{eq:v0-c5-1}:

\begin{equation}\label{eq:v0-c5-1}
\bm{v} = v_x\,\bm{i} + v_y\,\bm{j} + v_z\,\bm{k},
\qquad
v = \lVert\bm{v}\rVert = \sqrt{v_x^2 + v_y^2 + v_z^2}.
\end{equation}

Os vetores $\bm{i}$, $\bm{j}$, $\bm{k}$ são a base cartesiana: têm módulo
um e apontam nos eixos coordenados. Trocar de base muda as componentes,
não o vetor --- a velocidade de um rotor não depende de onde se coloca a
régua. A primeira habilidade operacional é decompor: dado um vetor de
módulo $v$ na direção do ângulo $\theta$ no plano, suas componentes são
$v_x = v\cos\theta$ e $v_y = v\,\mathrm{sen}\,\theta$.

\subsection{As duas multiplicações}

Entre vetores há duas multiplicações, e confundi-las é erro frequente. O
\textbf{produto escalar}, Equação~\eqref{eq:v0-c5-2}, devolve um número e
mede projeção: é o trabalho de uma força ao longo de um deslocamento e a
potência de um torque na velocidade angular.

\begin{equation}\label{eq:v0-c5-2}
\bm{a}\cdot\bm{b} = a\,b\cos\theta_{ab}
                  = a_x b_x + a_y b_y + a_z b_z.
\end{equation}

O \textbf{produto vetorial}, Equação~\eqref{eq:v0-c5-3}, devolve um vetor
perpendicular ao plano dos dois, com módulo igual à área do paralelogramo
que eles definem:

\begin{equation}\label{eq:v0-c5-3}
\lVert\bm{a}\times\bm{b}\rVert = a\,b\,\mathrm{sen}\,\theta_{ab},
\qquad
\bm{a}\times\bm{b} =
\begin{vmatrix}
\bm{i} & \bm{j} & \bm{k} \\
a_x & a_y & a_z \\
b_x & b_y & b_z
\end{vmatrix}.
\end{equation}

O momento de uma força, $\bm{\tau} = \bm{r}\times\bm{F}$, e a rotação
$\bm{v} = \bm{\omega}\times\bm{r}$ de um ponto a velocidade angular
$\bm{\omega}$ são produtos vetoriais: a direção perpendicular não é
detalhe, é a informação principal. A Figura~\ref{fig:v0-c5-1} resume as
duas operações no plano.

\begin{figure}[H]
\centering
\begin{tikzpicture}[font=\footnotesize, scale=1.5]
% --- (a) produto escalar: projeção ---
\begin{scope}[xshift=0cm]
\draw[->,thick] (0,0) -- (2.6,0) node[below=1pt] {$\bm{a}$};
\draw[->,thick] (0,0) -- (1.8,1.2) node[above=1pt] {$\bm{b}$};
\draw[dashed] (1.8,1.2) -- (1.8,0);
\draw[decorate,decoration={brace,amplitude=3pt},yshift=-2pt]
  (0,0) -- (1.8,0) node[midway,below=4pt] {$b\cos\theta$};
\draw (0.42,0) arc (0:33.7:0.42) node[midway,right=1pt] {$\theta$};
\node at (1.3,-0.75) {(a) escalar: projeção};
\end{scope}
% --- (b) produto vetorial: área ---
\begin{scope}[xshift=5.6cm]
\draw[->,thick] (0,0) -- (2.6,0) node[below=1pt] {$\bm{a}$};
\draw[->,thick] (0,0) -- (1.1,1.5) node[above=1pt] {$\bm{b}$};
\draw[black!45] (2.6,0) -- (3.7,1.5) -- (1.1,1.5);
\fill[black!12] (0,0) -- (2.6,0) -- (3.7,1.5) -- (1.1,1.5) -- cycle;
\draw (0.42,0) arc (0:53.7:0.42) node[midway,right=1pt] {$\theta$};
\node at (1.85,-0.75) {(b) vetorial: área $a\,b\,\mathrm{sen}\,\theta$};
\end{scope}
\end{tikzpicture}
\caption[Produto escalar como projeção e produto vetorial como área.]{%
As duas multiplicações entre vetores. (a) O produto escalar mede a
projeção de $\bm{b}$ sobre $\bm{a}$. (b) O módulo do produto vetorial é a
área do paralelogramo definido pelos dois vetores; seu sentido é dado pela
regra da mão direita, perpendicular ao plano.}
\label{fig:v0-c5-1}
\end{figure}

Duas identidades que a coleção usa sem aviso decorrem das
definições: o produto escalar comuta ($\bm{a}\cdot\bm{b} =
\bm{b}\cdot\bm{a}$), o produto vetorial anticomuta
($\bm{a}\times\bm{b} = -\,\bm{b}\times\bm{a}$), e um vetorial duplo
repetido é nulo ($\bm{a}\times\bm{a} = \bm{0}$).

\subsection{Matrizes e sistemas lineares}

Uma matriz é uma tabela retangular de números que representa uma
transformação linear: multiplicada por um vetor-coluna, devolve outro
vetor. O que a engenharia de máquinas precisa saber operar é pequeno: o
produto linha por coluna, a transposição (trocar linhas por colunas, com
$[\,\cdot\,]^{T}$), e a inversão --- a matriz $\bm{A}^{-1}$ tal que
$\bm{A}^{-1}\bm{A} = \bm{I}$, onde $\bm{I}$ é a identidade. Sistemas
lineares escrevem-se compactamente como $\bm{A}\bm{x} = \bm{b}$, e a
solução, quando existe, é $\bm{x} = \bm{A}^{-1}\bm{b}$.

O conteúdo prático está em saber \emph{quando} a solução existe e é única:
quando o determinante de $\bm{A}$ é não nulo, ou, em linguagem física,
quando as equações do sistema são independentes entre si. Um sistema com
duas equações que dizem a mesma coisa tem determinante nulo e infinitas
soluções; um sistema com equações contraditórias tem determinante nulo e
nenhuma solução. O Volume~II encontra exatamente esses casos nos balanços
de redes de ar e nos ajustes de parâmetros por mínimos quadrados, onde a
matriz dos coeficientes precisa ser invertida numericamente.

\subsection{Síntese: premissas e ligação com os volumes seguintes}

Este capítulo definiu vetor, produtos escalar e vetorial, matriz e
inversão. Os produtos das Equações~\eqref{eq:v0-c5-2}
e~\eqref{eq:v0-c5-3} reaparecem na mecânica do contínuo (tensão e
momento), na eletricidade (potência e campo) e na dinâmica rotacional do
Volume~I; a inversão de matrizes sustenta a resolução numérica do
Volume~II.


%=============================================================================
% CAPÍTULO 6 — CÁLCULO DIFERENCIAL
%=============================================================================
\section{Cálculo Diferencial}\label{sec:v0-derivadas}

\subsection{A derivada como taxa instantânea}

A pergunta que o cálculo diferencial responde é: com que rapidez uma
grandeza muda naquele instante? A velocidade média entre dois instantes é
$\Delta x / \Delta t$; a velocidade instantânea é o limite desse quociente
quando o intervalo encolhe, Equação~\eqref{eq:v0-c6-1}:

\begin{equation}\label{eq:v0-c6-1}
\frac{\mathrm{d}x}{\mathrm{d}t}
  = \lim_{\Delta t \to 0} \frac{\Delta x}{\Delta t}.
\end{equation}

Essa é a definição, e ela importa: toda derivada dos volumes seguintes é
um limite como este, e toda regra de derivação é um atalho para calculá-lo.
A notação $\mathrm{d}x/\mathrm{d}t$ lê-se ``taxa de variação de $x$ com
$t$''; quando a variável independente é o tempo, o ponto por cima
($\dot{x}$) abrevia a mesma coisa --- como na vazão $\dot{m}$ do
Volume~I.

\subsection{Regras de derivação}

As regras que a coleção usa estão na Tabela~\ref{tab:v0-c6-1}. Nenhuma
delas deve ser procurada em tabela durante a leitura dos volumes: todas
precisam estar disponíveis de imediato.

\begin{table}[H]
\centering
\caption[Regras de derivação usadas nesta coleção.]{%
Derivadas das funções elementares e regras operacionais.}
\label{tab:v0-c6-1}
\begin{tabularx}{\textwidth}{lX}
\toprule
\textbf{Função} & \textbf{Derivada} \\
\midrule
$x^n$ & $n\,x^{n-1}$ \\
$\mathrm{e}^{x}$ & $\mathrm{e}^{x}$ \\
$\ln x$ & $1/x$ \\
$\mathrm{sen}\,x$ & $\cos x$ \\
$\cos x$ & $-\,\mathrm{sen}\,x$ \\
\midrule
$f + g$ & $f' + g'$ \\
$f\,g$ & $f'g + f\,g'$ \\
$f(g(x))$ & $f'(g)\,g'(x)$ \quad (regra da cadeia) \\
$f/g$ & $(f'g - f\,g')/g^2$ \\
\bottomrule
\end{tabularx}
\end{table}

A regra da cadeia merece destaque porque é a mais usada e a mais
esquecida: ela diz que a taxa de uma grandeza composta é o produto das
taxas encadeadas. Quando o Volume~I escreve a aceleração de uma partícula
como $a = \mathrm{d}v/\mathrm{d}t = (\mathrm{d}v/\mathrm{d}x)(\mathrm{d}x/
\mathrm{d}t)$, é a regra da cadeia trabalhando.

\subsection{A derivada como inclinação --- e a linearização}

Geometricamente, a derivada num ponto é a inclinação da reta tangente à
curva naquele ponto, Figura~\ref{fig:v0-c6-1}. Essa leitura geométrica é
operacionalmente a mais importante: \emph{próximo de um ponto, toda função
diferenciável parece uma reta}. A Equação~\eqref{eq:v0-c6-2} é a
aproximação de primeira ordem, base de toda linearização da coleção:

\begin{equation}\label{eq:v0-c6-2}
f(x_0 + \Delta x) \approx f(x_0) + f'(x_0)\,\Delta x.
\end{equation}

\begin{figure}[H]
\centering
\begin{tikzpicture}[font=\footnotesize, scale=1.4]
\draw[->] (0,0) -- (6.3,0) node[below left=1pt] {$x$};
\draw[->] (0,0) -- (0,4.1) node[below left=1pt] {$f$};
% curva f(x) = 0.08 x^2 + 0.4 sen(2x) + 1 ... suave
\draw[thick, domain=0:6, samples=120]
  plot (\x, {0.09*\x*\x + 0.5*sin(\x r) + 0.8});
% ponto x0=2.6
\coordinate (P) at (2.6, {0.09*2.6*2.6 + 0.5*sin(2.6 r) + 0.8});
\fill (P) circle (1.4pt) node[above=1pt] {$f(x_0)$};
\draw[dashed] (2.6,0) -- (P);
\node[below] at (2.6,0) {$x_0$};
% tangente em x0: f'(x0) = 0.18*2.6 + 0.5*cos(2.6) = 0.468 + (-0.429) = 0.039
% usa inclinação real aprox 0.04; reta: y = f(x0) + 0.04 (x-2.6)
\draw[thick, dashed, domain=1.2:4.6]
  plot (\x, {0.09*2.6*2.6 + 0.5*sin(2.6 r) + 0.8 + 0.04*(\x-2.6)});
% secante até x0+dx
\coordinate (Q) at (4.3, {0.09*4.3*4.3 + 0.5*sin(4.3 r) + 0.8});
\draw[black!50] (P) -- (Q);
\draw[dashed] (4.3,0) -- (Q);
\node[below] at (4.3,0) {$x_0 + \Delta x$};
\node[black!50,right] at (3.6,2.85) {secante};
\node at (4.9,1.6) {tangente};
\node at (1.1,2.3) {$f(x)$};
\end{tikzpicture}
\caption[A derivada como inclinação da tangente.]{%
A derivada $f'(x_0)$ é a inclinação da tangente à curva em $x_0$. A
secante entre $x_0$ e $x_0+\Delta x$ aproxima a tangente; a aproximação se
torna exata quando $\Delta x \to 0$. Perto de $x_0$, a tangente substitui a
curva: é a linearização da Equação~\eqref{eq:v0-c6-2}.}
\label{fig:v0-c6-1}
\end{figure}

Todo o estudo de estabilidade, de sensibilidade de parâmetros e de
pequenos sinais nos volumes seguintes é a Equação~\eqref{eq:v0-c6-2}
aplicada a modelos físicos: o Volume~II lineariza o comportamento do
compressor em torno de um ponto de operação para projetar o controle.

\subsection{Derivadas parciais e o gradiente}

Uma função de duas variáveis, $f(x,y)$, tem duas taxas de variação: a que
resulta de mexer só em $x$ (com $y$ fixo) e a que resulta de mexer só em
$y$. São as derivadas parciais, Equação~\eqref{eq:v0-c6-3}, escritas com o
sinal $\partial$ em vez de $\mathrm{d}$ para lembrar que a outra variável
ficou parada:

\begin{equation}\label{eq:v0-c6-3}
\frac{\partial f}{\partial x}
  = \lim_{\Delta x \to 0}
    \frac{f(x+\Delta x, y) - f(x,y)}{\Delta x}.
\end{equation}

O vetor das derivadas parciais é o gradiente,
$\bm{\nabla} f = (\partial f/\partial x,\ \partial f/\partial y,\ 
\partial f/\partial z)$: ele aponta na direção de maior crescimento de
$f$ e seu módulo mede a inclinação máxima. O gradiente de temperatura
dirige o fluxo de calor na lei de Fourier; o gradiente de pressão dirige o
escoamento; o gradiente de uma função de custo dirige os ajustes de
projeto. A regra da cadeia estende-se naturalmente: se $f$ depende de $x$
e de $y$, e ambas dependem de $t$, então

\begin{equation}\label{eq:v0-c6-4}
\frac{\mathrm{d}f}{\mathrm{d}t}
  = \frac{\partial f}{\partial x}\frac{\mathrm{d}x}{\mathrm{d}t}
  + \frac{\partial f}{\partial y}\frac{\mathrm{d}y}{\mathrm{d}t}.
\end{equation}

A Equação~\eqref{eq:v0-c6-4} é a versão contínua da contabilidade de
variações usada nos balanços do Volume~I: a variação total é a soma das
variações parciais vezes as taxas de cada caminho.

\subsection{Máximos, mínimos e pontos de equilíbrio}

Onde uma função suave atinge máximo ou mínimo, a tangente é horizontal:
$f'(x) = 0$. Esses pontos \emph{estacionários} são candidatos a ótimos, e
a segunda derivada decide: $f'' > 0$ indica mínimo, $f'' < 0$ indica
máximo, $f'' = 0$ pede investigação adicional. Em funções de duas
variáveis, o ponto estacionário exige $\partial f/\partial x = 0$ e
$\partial f/\partial y = 0$ simultaneamente. A busca de máximo
rendimento, de mínimo consumo específico e de ponto de operação ótimo,
que percorre o Volume~II do início ao fim, é esta seção em escala
industrial.

\subsection{Síntese: premissas e ligação com os volumes seguintes}

A derivada foi definida como limite de taxa média, e todas as suas regras
são consequência da definição. As três leituras --- taxa instantânea
(Equação~\eqref{eq:v0-c6-1}), inclinação e linearização
(Equação~\eqref{eq:v0-c6-2}) e variação parcial
(Equação~\eqref{eq:v0-c6-3}) --- cobrem todo o uso que os Volumes~I e~II
fazem do cálculo diferencial: leis de taxa, sensibilidades e otimização.


%=============================================================================
% CAPÍTULO 7 — CÁLCULO INTEGRAL
%=============================================================================
\section{Cálculo Integral}\label{sec:v0-integrais}

\subsection{A integral como acumulação}

Se a derivada pergunta ``com que rapidez?'', a integral pergunta ``quanto
acumulou?''. Uma vazão $\dot{m}$ constante durante $\Delta t$ acumula
$\dot{m}\,\Delta t$ de massa; uma vazão variável acumula a soma de muitos
pedacinhos, Equação~\eqref{eq:v0-c7-1}:

\begin{equation}\label{eq:v0-c7-1}
m = \int_{t_1}^{t_2} \dot{m}(t)\ \mathrm{d}t.
\end{equation}

O símbolo $\int$ é um ``S'' esticado, de soma: a integral soma produtos
$\dot{m}(t)\,\mathrm{d}t$, cada um sendo quantidade por vez vezes um
pedacinho de tempo. A interpretação é sempre essa: a integral acumula a
taxa. Todo balanço do Volume~I --- massa, energia, quantidade de
movimento --- é uma Equação~\eqref{eq:v0-c7-1} vestida de princípio de
conservação.

\subsection{O teorema fundamental}

Derivar e integrar são operações inversas, e isso é o teorema fundamental
do cálculo, Equação~\eqref{eq:v0-c7-2}: a integral da taxa devolve a
variação acumulada entre os extremos.

\begin{equation}\label{eq:v0-c7-2}
\int_{a}^{b} f'(x)\ \mathrm{d}x = f(b) - f(a).
\end{equation}

Duas consequências operacionais seguem. A primeira: para avaliar uma
integral definida, basta encontrar uma função cuja derivada seja o
integrando (uma \emph{primitiva}) e subtrair seus valores nos extremos.
A segunda: a integral definida é um \emph{número}, e a primitiva é uma
\emph{função}; confundir as duas coisas é a fonte clássica de erro de
sinal e de constante.

\subsection{A integral como área --- e o valor médio}

Geometricamente, a integral definida é a área sob a curva entre os
extremos, Figura~\ref{fig:v0-c7-1}, com áreas abaixo do eixo contadas
negativo. A leitura de área resolve problemas que a primitiva resolve com
dificuldade: a energia de um ciclo no diagrama $p$--$V$ é a área
fechada pela curva, e trabalho líquido de um ciclo fechado é zero quando a
curva não envolve área.

\begin{figure}[H]
\centering
\begin{tikzpicture}[font=\footnotesize, scale=1.4]
\draw[->] (0,0) -- (6.3,0) node[below left=1pt] {$x$};
\draw[->] (0,0) -- (0,3.6) node[below left=1pt] {$f$};
\draw[thick, domain=0:6, samples=120]
  plot (\x, {2.4*exp(-(\x-3)^2/2.2) + 0.35});
% área sob a curva entre x=1 e x=4.6
\fill[black!15, domain=1:4.6, samples=80]
  plot (\x, {2.4*exp(-(\x-3)^2/2.2) + 0.35}) -- (4.6,0) -- (1,0) -- cycle;
\draw[thick, domain=1:4.6, samples=80]
  plot (\x, {2.4*exp(-(\x-3)^2/2.2) + 0.35});
\draw[dashed] (1,0) -- (1,{2.4*exp(-4/2.2)+0.35});
\draw[dashed] (4.6,0) -- (4.6,{2.4*exp(-(1.6)^2/2.2)+0.35});
\node[below] at (1,0) {$a$};
\node[below] at (4.6,0) {$b$};
\node at (2.8,1.0) {$\displaystyle\int_a^b f(x)\,\mathrm{d}x$};
\end{tikzpicture}
\caption[A integral definida como área sob a curva.]{%
A integral definida de $f$ entre $a$ e $b$ é a área sob a curva nesse
intervalo. A região sombreada representa a acumulação da Equação
\eqref{eq:v0-c7-1}: cada fatia vertical de largura $\mathrm{d}x$ tem área
$f(x)\,\mathrm{d}x$, e a integral soma todas.}
\label{fig:v0-c7-1}
\end{figure}

Uma aplicação direta define o \textbf{valor médio} de uma função num
intervalo, Equação~\eqref{eq:v0-c7-3}:

\begin{equation}\label{eq:v0-c7-3}
\bar{f} = \frac{1}{b-a}\int_{a}^{b} f(x)\ \mathrm{d}x.
\end{equation}

O valor médio é a altura do retângulo de mesma área que a curva. A
diferença média logarítmica de temperatura dos trocadores de calor e a
potência média de um ciclo periódico, ambas do Volume~I, são
Equação~\eqref{eq:v0-c7-3} com nomes próprios.

\subsection{Primitivas das funções elementares}

A Tabela~\ref{tab:v0-c7-1} lista as primitivas que a coleção usa. Como
verificação de segurança, derive o resultado: se a derivada volta ao
integrando, a primitiva está certa.

\begin{table}[H]
\centering
\caption[Primitivas das funções elementares.]{%
Primitivas imediatas das funções elementares, com $C$ constante arbitrária.}
\label{tab:v0-c7-1}
\begin{tabularx}{\textwidth}{lX}
\toprule
\textbf{Função} & \textbf{Primitiva} \\
\midrule
$x^n$ ($n \neq -1$) & $x^{n+1}/(n+1) + C$ \\
$1/x$ & $\ln|x| + C$ \\
$\mathrm{e}^{x}$ & $\mathrm{e}^{x} + C$ \\
$\mathrm{sen}\,x$ & $-\cos x + C$ \\
$\cos x$ & $\mathrm{sen}\,x + C$ \\
$1/(1+x^2)$ & $\operatorname{arctg} x + C$ \\
\bottomrule
\end{tabularx}
\end{table}

Quando a primitiva não é imediata, restam três caminhos, em ordem de
preferência: substituição de variável (que desfaz a regra da cadeia),
integração por partes (que desfaz a regra do produto) e tabelas ou
software para as demais. Para a engenharia de máquinas, o essencial é
reconhecer \emph{qual} caminho se aplica, e a coleção só exige isso.

\subsection{Integrais múltiplas}

Uma integral dupla $\iint_A f(x,y)\,\mathrm{d}A$ soma contribuições sobre
uma região do plano --- a massa de uma placa, por exemplo, soma
$\rho(x,y)\,\mathrm{d}A$ sobre toda a área. A avaliação é sempre a mesma:
integrar duas vezes em sequência, cada vez tratando a outra variável como
constante. A mesma lógica se estende a três dimensões. No Volume~I, os
balanços integrais sobre volumes de controle são exatamente integrais
triplas disfarçadas de física; o Capítulo~4 escreve o teorema do
transporte que as organiza.

\subsection{Síntese: premissas e ligação com os volumes seguintes}

A integral acumula taxa (Equação~\eqref{eq:v0-c7-1}), a primitiva inverte a
derivação (Equação~\eqref{eq:v0-c7-2}) e o valor médio resume o
acumulado (Equação~\eqref{eq:v0-c7-3}). Com isso, o leitor está pronto
para os balanços integrais do Volume~I e para a acumulação de energia,
massa e carga que percorre o Volume~II.


%=============================================================================
% CAPÍTULO 8 — EQUAÇÕES DIFERENCIAIS ORDINÁRIAS
%=============================================================================
\section{Equações Diferenciais Ordinárias}\label{sec:v0-edo}

\subsection{O que é uma equação diferencial}

Uma equação diferencial ordinária (EDO) relaciona uma função às suas
próprias taxas de variação. A declaração ``a temperatura do reservatório
cai a uma taxa proporcional à diferença entre ela e a ambiente'' traduz-se
na Equação~\eqref{eq:v0-c8-1}:

\begin{equation}\label{eq:v0-c8-1}
\frac{\mathrm{d}T}{\mathrm{d}t} = -k\,(T - T_{\infty}),
\end{equation}

com $k > 0$ constante do processo e $T_{\infty}$ a temperatura ambiente.
A solução não é um número, é uma função $T(t)$ que satisfaz a equação para
todo $t$. A \emph{ordem} da EDO é a maior derivada presente; a
Equação~\eqref{eq:v0-c8-1} é de primeira ordem. Quase toda dinâmica do
Volume~II --- resposta de pressão, de temperatura, de rotação --- é uma
EDO de primeira ou segunda ordem com nomes próprios.

\subsection{Condições iniciais e o problema completo}

Uma EDO admite infinitas soluções: a Equação~\eqref{eq:v0-c8-1} é satisfeita
por qualquer decaimento exponencial de amplitude livre. Fixar o problema
exige \emph{condições iniciais}: o valor da função (e das derivadas, se a
ordem for maior que um) num instante de referência. O conjunto EDO mais
condição inicial chama-se \emph{problema de valor inicial}, e sua solução
é única sempre que a equação for suficientemente bem comportada --- o que
todas as equações desta coleção são. Saber contar as condições exigidas é
metade da modelagem: primeira ordem pede uma condição, segunda ordem pede
duas.

\subsection{Primeira ordem: separáveis e lineares}

A Equação~\eqref{eq:v0-c8-1} é do tipo \emph{separável}: as variáveis se
separam, com tudo que depende de $T$ de um lado e tudo que depende de $t$
do outro, Equação~\eqref{eq:v0-c8-2}:

\begin{equation}\label{eq:v0-c8-2}
\frac{\mathrm{d}T}{T - T_{\infty}} = -k\ \mathrm{d}t
\quad\Longrightarrow\quad
\ln\frac{T - T_{\infty}}{T_0 - T_{\infty}} = -k\,t,
\end{equation}

onde $T_0 = T(0)$. Isolando $T(t)$, obtém-se o decaimento exponencial,
Equação~\eqref{eq:v0-c8-3}, com a constante de tempo $\tau = 1/k$:

\begin{equation}\label{eq:v0-c8-3}
T(t) = T_{\infty} + (T_0 - T_{\infty})\,\mathrm{e}^{-t/\tau}.
\end{equation}

A constante de tempo tem leitura direta: a cada intervalo $\tau$, a
diferença $T - T_{\infty}$ cai para $37\%$ ($\mathrm{e}^{-1}$) do valor
anterior; em $3\tau$, restam $5\%$. Esse argumento aparece em todos os
sistemas de primeira ordem da coleção, do resfriamento de motor à
resposta de sensor.

A segunda família de primeira ordem é a \emph{linear}:
$\dot{x} + a(t)\,x = b(t)$, que inclui forçamento externo $b(t)$. A
estratégia de solução (fator integrante) não precisa ser reproduzida aqui;
o que importa é a estrutura da solução: resposta = transiente (decai com
o tempo) + regime permanente (acompanha o forçamento). Essa decomposição
organiza a análise de resposta do Volume~II.

\subsection{Segunda ordem: oscilações}

A EDO de segunda ordem com coeficientes constantes, Equação
\eqref{eq:v0-c8-4}, governa massa--mola--amortecedor, circuitos $RLC$ e
todo oscilador natural de máquina:

\begin{equation}\label{eq:v0-c8-4}
\ddot{x} + 2\zeta\omega_n\,\dot{x} + \omega_n^2\,x = f(t).
\end{equation}

O parâmetro $\omega_n$ é a \emph{frequência natural} (a que o sistema
oscilaria sem amortecimento) e $\zeta$ é o \emph{fator de amortecimento},
que classifica o comportamento, Figura~\ref{fig:v0-c8-1}:

\begin{itemize}
\item $\zeta < 1$ --- subamortecido: oscila com amplitude decaindo;
\item $\zeta = 1$ --- criticamente amortecido: retorna ao repouso no menor
  tempo sem cruzar a posição de equilíbrio;
\item $\zeta > 1$ --- superamortecido: retorna sem oscilar, lentamente.
\end{itemize}

\begin{figure}[H]
\centering
\begin{tikzpicture}[font=\footnotesize, scale=1.15]
\begin{axis}[
  width=11.5cm, height=6.2cm,
  xlabel={$t$ (tempo adimensional $\omega_n t$)},
  ylabel={$x$ (deslocamento)},
  xmin=0, xmax=12, ymin=-0.4, ymax=1.25,
  axis lines=left,
  legend style={font=\footnotesize, at={(0.97,0.03)}, anchor=south east},
  xtick=\empty, ytick={0,1},
]
% subamortecido zeta=0.15: x = exp(-0.15 t) cos(0.989 t)
\addplot[thick, domain=0:12, samples=240]
  {1 - exp(-0.15*x)*(cos(deg(0.9887*x)) + 0.1517*sin(deg(0.9887*x)))};
\addlegendentry{$\zeta = 0{,}15$ (subamortecido)}
% criticamente amortecido: x = (1+t) exp(-t)
\addplot[thick, dashed, domain=0:12, samples=120]
  {1 - (1+x)*exp(-x)};
\addlegendentry{$\zeta = 1$ (crítico)}
% superamortecido zeta=2: raízes -0.268 e -3.732; x = (r2 e^{r1 t} - r1 e^{r2 t})/(r2-r1)
\addplot[thick, dotted, domain=0:12, samples=120]
  {1 + (-3.732*exp(-0.268*x) + 0.268*exp(-3.732*x))/(-0.268+3.732)};
\addlegendentry{$\zeta = 2$ (superamortecido)}
\end{axis}
\end{tikzpicture}
\caption[Resposta de sistemas de segunda ordem por regime de amortecimento.]{%
Resposta ao degrau de um sistema de segunda ordem, Equação
\eqref{eq:v0-c8-4}, para três fatores de amortecimento. O subamortecido
oscila antes de assentar; o crítico assenta sem oscilar no menor tempo; o
superamortecido assenta sem oscilar, mais lentamente.}
\label{fig:v0-c8-1}
\end{figure}

O Volume~I deduz a Equação~\eqref{eq:v0-c8-4} três vezes --- mecânica,
eletricidade e vibrações --- e mostra que é a mesma equação com roupagens
diferentes; o Volume~II a usa para avaliar se o comportamento dinâmico de
um compressor está próximo de um regime perigoso de oscilação.

\subsection{Solução numérica: o passo de Euler}

Quando a EDO não admite solução fechada --- e as EDOs com coeficientes
variáveis ou não lineares, como as do Volume~II, em geral não admitem ---
recorre-se à integração numérica. O método de Euler, o mais simples,
traduz a definição de derivada direto em algoritmo: se
$\dot{x} = g(x,t)$, então a variação num passo pequeno $\Delta t$ é
aproximadamente $g\,\Delta t$, Equação~\eqref{eq:v0-c8-5}:

\begin{equation}\label{eq:v0-c8-5}
x_{n+1} = x_n + g(x_n, t_n)\ \Delta t.
\end{equation}

A Equação~\eqref{eq:v0-c8-5} é a própria definição da derivada com o
limite suspenso; por isso o erro de cada passo é proporcional a
$(\Delta t)^2$, e o erro acumulado, a $\Delta t$. Refinar o passo
diminui o erro previsivelmente --- e essa previsibilidade é o que permite
confiar na simulação. O Volume~II integra suas EDOs com métodos mais
refinados que este, mas todos descendentes diretos da
Equação~\eqref{eq:v0-c8-5}.

\subsection{Síntese: premissas e ligação com os volumes seguintes}

Uma EDO postula a forma da taxa; a condição inicial fixa a amplitude; a
solução fecha o problema. As famílias aqui apresentadas --- primeira ordem
separável e linear, segunda ordem com coeficientes constantes --- cobrem
toda a dinâmica dos volumes seguintes. A constante de tempo
(Equação~\eqref{eq:v0-c8-3}) e o fator de amortecimento
(Equação~\eqref{eq:v0-c8-4}) serão vocabulário corrente a partir daqui.


%=============================================================================
% CAPÍTULO 9 — SÉRIES E APROXIMAÇÕES
%=============================================================================
\section{Séries e Aproximações}\label{sec:v0-series}

\subsection{Por que aproximar}

Os modelos exatos são raros; os úteis são aproximados. A pergunta da
engenharia nunca é ``a aproximação é perfeita?'', e sim ``o erro da
aproximação é pequeno \emph{para o meu propósito}?''. Este capítulo fornece
as duas ferramentas para respondê-la: a série de Taylor, que gera
aproximações, e a notação de ordem, que declara o erro.

\subsection{A série de Taylor}

Uma função suave pode ser reconstruída perto de um ponto a partir de seus
valores e de todas as suas derivadas naquele ponto, Equação
\eqref{eq:v0-c9-1}:

\begin{equation}\label{eq:v0-c9-1}
f(x_0 + h) = f(x_0) + f'(x_0)\,h
  + \frac{f''(x_0)}{2!}\,h^2
  + \frac{f'''(x_0)}{3!}\,h^3 + \dots
\end{equation}

O polinômio truncado na primeira ordem é a linearização da
Equação~\eqref{eq:v0-c6-2}; cada termo a mais melhora a cópia numa
vizinhança maior. O fatorial nos denominadores faz a série convergir
rapidamente para $h$ pequeno: o termo de ordem $n$ decai com $h^n$, de
modo que dobrar a precisão da aproximação custa cada vez menos.

A Figura~\ref{fig:v0-c9-1} mostra a aproximação de $\mathrm{e}^{x}$ em
torno de $x_0 = 0$: a reta (ordem 1) serve numa vizinhança estreita, a
parábola (ordem 2) alarga-a, e a cúbica (ordem 3) praticamente coincide
com a exponencial no intervalo desenhado.

\begin{figure}[H]
\centering
\begin{tikzpicture}[font=\footnotesize]
\begin{axis}[
  width=11.5cm, height=6.6cm,
  xlabel={$x$}, ylabel={$f(x)$},
  xmin=-2.2, xmax=2.2, ymin=-0.5, ymax=4.5,
  axis lines=left,
  legend style={font=\footnotesize, at={(0.03,0.97)}, anchor=north west},
  ytick={0,1,2,3,4},
]
\addplot[thick, domain=-2.2:1.5, samples=120] {exp(x)};
\addlegendentry{$\mathrm{e}^{x}$}
\addplot[thick, dashed, domain=-2.2:2.2] {1+x};
\addlegendentry{ordem 1: $1+x$}
\addplot[thick, dotted, domain=-2.2:2.2] {1+x+x^2/2};
\addlegendentry{ordem 2: $1+x+x^2/2$}
\addplot[thin, domain=-2.2:2.2] {1+x+x^2/2+x^3/6};
\addlegendentry{ordem 3: $+x^3/6$}
\end{axis}
\end{tikzpicture}
\caption[Série de Taylor da exponencial em torno de $x_0=0$.]{%
A série de Taylor da Equação~\eqref{eq:v0-c9-1} aplicada a
$\mathrm{e}^{x}$ em torno de $x_0 = 0$. Cada ordem acrescentada melhora a
aproximação numa vizinhança maior; perto de $x_0$, a reta já descreve a
função com erro pequeno.}
\label{fig:v0-c9-1}
\end{figure}

As expansões que a coleção usa sem aviso decorrem da
Equação~\eqref{eq:v0-c9-1}, válidas para $|\varepsilon| \ll 1$,
Equação~\eqref{eq:v0-c9-2}:

\begin{equation}\label{eq:v0-c9-2}
(1+\varepsilon)^n \approx 1 + n\,\varepsilon,
\qquad
\frac{1}{1+\varepsilon} \approx 1 - \varepsilon,
\qquad
\mathrm{e}^{\varepsilon} \approx 1 + \varepsilon,
\qquad
\ln(1+\varepsilon) \approx \varepsilon.
\end{equation}

A primeira delas é o binômio de Newton truncado; todas dizem o mesmo:
funções suaves parecem retas de perto, e a declaração
``$|\varepsilon| \ll 1$'' é a hipótese que dá sentido à aproximação.

\subsection{Erro e notação de ordem}

A notação $f(h) = O(h^n)$ lê-se ``$f$ é da ordem de $h^n$'' e significa
que $f$ decai pelo menos tão rápido quanto $h^n$ quando $h \to 0$:
existem constantes $C$ e $h_0$ tais que $|f(h)| \le C\,h^n$ para
$0 < h < h_0$. A constante $C$ não interessa no momento --- interessa o
expoente, porque ele diz como o erro melhora quando se refina a
discretização. O método de Euler da Equação~\eqref{eq:v0-c8-5} tem erro
global $O(\Delta t)$: diminuir o passo pela metade divide o erro pela
metade, e essa previsão é verificável numericamente.

Declarar a ordem do erro é o mínimo de honestidade de uma aproximação.
Quando o Volume~II lineariza um modelo, a pergunta ``qual é o $h$ da
linearização?'' --- isto é, quão longe do ponto de operação a aproximação
ainda vale --- tem resposta numérica, e a notação $O(\cdot)$ é a
linguagem para formular essa resposta.

\subsection{Síntese: premissas e ligação com os volumes seguintes}

A série de Taylor gera aproximações locais; a notação de ordem declara seu
erro. Juntas, autorizam a operação mais frequente desta coleção:
substituir um modelo complicado por um simples \emph{numa vizinhança
declarada}, sabendo quanto se perde. Todas as linearizações, todas as
sensibilidades e todos os controladores dos volumes seguintes usam essa
autorização.


%=============================================================================
% CAPÍTULO 10 — SÍNTESE E PONTE COM OS VOLUMES I E II
%=============================================================================
\section{Síntese e Ponte com os Volumes I e II}\label{sec:v0-sintese}

\subsection{As peças reunidas}

O Volume~0 não apresentou física nenhuma: apresentou a linguagem e as
operações com que a física será escrita. A Tabela~\ref{tab:v0-c10-1}
fecha o ciclo começado na Tabela~\ref{tab:v0-c1-1}: cada ferramenta, o
capítulo que a ensina e o primeiro uso que os volumes seguintes fazem dela.

\begin{table}[H]
\centering
\caption[Ferramentas do Volume 0 e seu primeiro uso nos volumes seguintes.]{%
Mapa de saída do Volume 0: de cada ferramenta ao seu primeiro uso na
coleção.}
\label{tab:v0-c10-1}
\begin{tabularx}{\textwidth}{lXl}
\toprule
\textbf{Cap.} & \textbf{Ferramenta} & \textbf{Primeiro uso} \\
\midrule
\ref{sec:v0-notacao} & Notação & Símbolos das leis de taxa, Cap. 2 do Vol.~I \\
\ref{sec:v0-unidades} & Unidades e dimensão & Verificação de toda equação do Vol.~I \\
\ref{sec:v0-funcoes} & Funções & Leis constitutivas, Caps. 2--3 do Vol.~I \\
\ref{sec:v0-vetores} & Vetores & Tensão de Cauchy, Cap. 2 do Vol.~I \\
\ref{sec:v0-derivadas} & Derivadas & Leis de taxa, Caps. 2--5 do Vol.~I \\
\ref{sec:v0-integrais} & Integrais & Balanços integrais, Cap. 4 do Vol.~I \\
\ref{sec:v0-edo} & Equações diferenciais & Sistemas dinâmicos, Vol.~II \\
\ref{sec:v0-series} & Séries e ordem & Linearizações, Vols.~I e~II \\
\bottomrule
\end{tabularx}
\end{table}

\subsection{Como o Volume 0 se relaciona com os demais}

A relação é de pré-requisito, não de resumo: o Volume~I pressupõe este
volume da mesma forma que prescreve seus próprios capítulos. Quem leu
aqui não encontra surpresa simbólica adiante; quem pulou aqui e tropeçar
num símbolo ou numa operação encontra neste volume o lugar de onde ele
veio. A convenção de notação do Capítulo~\ref{sec:v0-notacao} vale sem
exceção: se algum símbolo parecer usado de outro modo, o uso local está
errado, e não a convenção.

\subsection{Próximos passos}

Com este volume concluído, o percurso natural segue para o Volume~I, que
aplica estas ferramentas à mecânica do contínuo, aos termofluidos e aos
sistemas de medição --- e destes, ao Volume~II, que as particulariza no
compressor de parafuso offshore. Quem chegou aqui sem formação prévia
pode precisar de um percurso mais longo; Nussenzveig~\cite{nussenzveig}
fornece a física básica que antecede este volume, e Guidorizzi
\cite{guidorizzi2013} o cálculo em profundidade. As referências de classe
mundial para o mesmo percurso são Feynman~\cite{feynman1963} na física e
Apostol~\cite{apostol1967} e Spivak~\cite{spivak2008} no cálculo. O
caminho da coleção, porém, está completo: o que será usado, já foi
ensinado.


%=============================================================================
% BIBLIOGRAFIA
%=============================================================================
\newpage
\begin{thebibliography}{99}

\bibitem{nussenzveig} Nussenzveig, H. M. \textit{Curso de Física Básica}. Vol.~1: \textit{Mecânica}. 5ª ed. São Paulo: Blucher, 2013; Vol.~2: \textit{Fluidos, Oscilações e Ondas, Calor}. 5ª ed. São Paulo: Blucher, 2014.

\bibitem{guidorizzi2013} Guidorizzi, H. L. \textit{Um Curso de Cálculo}. Vol.~1. 5ª ed. Rio de Janeiro: LTC, 2013.

\bibitem{feynman1963} Feynman, R. P., Leighton, R. B., Sands, M. \textit{The Feynman Lectures on Physics}. Vol.~1: \textit{Mainly Mechanics, Radiation, and Heat}. Addison-Wesley, 1963.

\bibitem{apostol1967} Apostol, T. M. \textit{Calculus}. Vol.~1: \textit{One-Variable Calculus, with an Introduction to Linear Algebra}. 2nd ed. Wiley, 1967.

\bibitem{spivak2008} Spivak, M. \textit{Calculus}. 4th ed. Houston: Publish or Perish, 2008.

\end{thebibliography}

\end{document}
